\section{The Constitutional Mapping}
\label{sec:mapping}

We establish a strict correspondence between Montesquieu's trisection of
institutional power~\cite{montesquieu1748} and the architectural layers
of the BTV stack. The correspondence is not analogical but
\emph{structural}: each political requirement maps to a formal property
that is either enforced at compile-time or cryptographically guaranteed
at runtime. We first define constitutional invariants, then prove
independence for each branch.

\subsection{Constitutional Invariants}
\label{subsec:invariants}

\begin{definition}[Algorithmic Institution]
\label{def:institution}
An \emph{algorithmic institution} $\mathcal{I}$ is a tuple
$\langle \mathcal{L}, \mathcal{E}, \mathcal{J}, \Sigma \rangle$ where:
\begin{itemize}
  \item $\mathcal{L}$ is the \emph{legislative layer}: a set of type
  constraints $\tau \in \mathcal{T}$ encoding compliance rules;
  \item $\mathcal{E}$ is the \emph{executive layer}: a program
  $P : \mathcal{T} \to \Sigma$ that produces signed state transitions
  under the constraints of $\mathcal{L}$;
  \item $\mathcal{J}$ is the \emph{judicial layer}: a verification
  function $V : \Sigma \to \{0,1\}$ that audits state integrity without
  querying $\mathcal{E}$;
  \item $\Sigma$ is the \emph{global state}: an append-only,
  content-addressed log of signed state transitions.
\end{itemize}
\end{definition}

\begin{definition}[Constitutional Invariant]
\label{def:invariant}
An algorithmic institution $\mathcal{I}$ satisfies the
\emph{Constitutional Invariant} if and only if the following three
conditions hold simultaneously:
\begin{align}
  P_{\mathcal{L}} \cap P_{\mathcal{E}} &= \emptyset \label{eq:leg-exec} \\
  P_{\mathcal{E}} \cap P_{\mathcal{J}} &= \emptyset \label{eq:exec-jud} \\
  P_{\mathcal{J}} \cap P_{\mathcal{L}} &= \emptyset \label{eq:jud-leg}
\end{align}
where $P_{X}$ denotes the set of capabilities (read, write, delete,
redefine) held exclusively by layer $X$ over the global state $\Sigma$.
\end{definition}

The disjointness conditions~(\ref{eq:leg-exec}--\ref{eq:jud-leg})
formalize Montesquieu's core requirement: no single actor may
simultaneously define the rules \emph{and} execute them \emph{and}
judge compliance. We now show that the BTV architecture satisfies
all three conditions.

\subsection{The Constitutional Mapping Table}
\label{subsec:table}

Table~\ref{tab:mapping} establishes the correspondence between
political theory and BTV architecture. Each row names the political
concept, its BTV implementation, the formal property that discharges
it, and the paper in this series that provides the proof or
construction.

\begin{table*}[t]
\centering
\caption{Constitutional Correspondence: Montesquieu $\leftrightarrow$ BTV}
\label{tab:mapping}
\begin{tabular}{p{2.8cm} p{3.5cm} p{4.5cm} p{2.8cm} p{1.5cm}}
\toprule
\textbf{Political Concept} &
\textbf{BTV Layer} &
\textbf{Formal Guarantee} &
\textbf{Mechanism} &
\textbf{Ref.} \\
\midrule
\textbf{Legislative Power}\ (defines the law) &
Constitutional Type System $\mathcal{L}$ &
$\forall P \in \mathcal{E}:\; \Gamma \vdash P : \tau \Rightarrow \tau \in \mathcal{L}$\ (type-checks only under declared rules) &
Linear resource types; compile-time rejection &
\cite{btv2026} \\
\addlinespace
\textbf{Executive Power}\ (runs the state) &
Operator $\mathcal{E}$ &
$\forall \sigma \in \Sigma:\; \nexists\ \mathtt{delete}(\sigma)$ by $\mathcal{E}$\ (append-only; evidence cannot be suppressed) &
Merkle append-only log; non-repudiable receipts &
\cite{btv-t-repo} \\
\addlinespace
\textbf{Judicial Power}\ (verifies the state) &
Independent Auditor $\mathcal{J}$ &
$V(\sigma) \in \{0,1\}$ without querying $\mathcal{E}$, and\ $\forall i:\; \neg\mathsf{Reveal}(u_i)$\ (audits without operator permission; privacy-preserving) &
ZK inclusion proofs; public Merkle root &
\cite{btv-ar2026} \\
\addlinespace
\textbf{Separation}\ (powers are disjoint) &
Capability partition $P_{\mathcal{L}} \cap P_{\mathcal{E}} \cap P_{\mathcal{J}} = \emptyset$ &
$\mathcal{E}$ cannot redefine $\mathcal{L}$ types;\ $\mathcal{E}$ cannot block $\mathcal{J}$ verification;\ $\mathcal{J}$ cannot override $\mathcal{L}$ rules &
Type immutability; public log; ZK soundness &
This paper \\
\addlinespace
\textbf{Checks \& Balances}\ (each constrains the others) &
Composition guarantee &
$\mathcal{J}$ can invalidate $\mathcal{E}$ output;\ $\mathcal{L}$ rejects $\mathcal{E}$ programs at compile-time;\ $\mathcal{E}$ cannot forge $\mathcal{J}$-valid proofs &
Type system + ZK soundness + append-only &
This paper \\
\bottomrule
\end{tabular}
\end{table*}

\subsection{Legislative Independence}
\label{subsec:leg-independence}

\begin{theorem}[Legislative Independence]
\label{thm:leg}
In a BTV institution $\mathcal{I}$, the executive layer $\mathcal{E}$
cannot modify the constitutional type constraints $\mathcal{L}$ at
runtime.
\end{theorem}

\begin{proof}
By the linear type discipline of Paper~1~\cite{btv2026}, every program
$P \in \mathcal{E}$ is type-checked at compile time against a fixed
context $\Gamma$ derived from $\mathcal{L}$. The type context $\Gamma$
is an immutable artifact of the build process: it is not a runtime
parameter, not a database entry, and not accessible via any API exposed
to the operator. Formally, $\Gamma$ is a function of the source
dependencies $\mathcal{D}$ at compile time:
\[
  \Gamma = \Gamma(\mathcal{D}_{t_0})
\]
where $t_0$ is the compilation timestamp. Any change to $\mathcal{L}$
requires recompilation, which produces a new binary with a distinct
cryptographic hash. The production binary's hash is published to
$\Sigma$ at deployment (Paper~2~\cite{btv-t-repo}), so any substitution
of a differently-typed binary is detectable by $\mathcal{J}$. Therefore
$P_{\mathcal{E}} \cap P_{\mathcal{L}} = \emptyset$ at runtime.
\end{proof}

\subsection{Judicial Independence}
\label{subsec:jud-independence}

\begin{theorem}[Judicial Independence]
\label{thm:jud}
In a BTV institution $\mathcal{I}$, the auditor $\mathcal{J}$ can
verify the integrity of any state transition $\sigma \in \Sigma$
without querying the operator $\mathcal{E}$ and without learning the
private attributes of any individual decision subject.
\end{theorem}

\begin{proof}
The global state $\Sigma$ is a public, append-only Merkle tree
(Paper~2~\cite{btv-t-repo}). The Merkle root $r$ is published
periodically to a public bulletin. To audit inclusion of transition
$\sigma_i$, the auditor constructs or receives a ZK inclusion proof
$\pi_i$ (Paper~3~\cite{btv-ar2026}) such that:
\[
  \mathsf{Verify}(r,\, \pi_i,\, \mathsf{pub}(\sigma_i)) = 1
\]
where $\mathsf{pub}(\sigma_i)$ is the public component of the transition
(decision type, timestamp, compliance tag) and the private attributes
(individual identity, sensitive inputs) remain hidden behind the ZK
commitment. The verification call does not involve $\mathcal{E}$: it
requires only the public root $r$ and the proof $\pi_i$. By the
soundness property of the ZK proof system~\cite{btv-ar2026}, a forged
proof for a non-member transition is accepted with probability at most
$\epsilon_{\mathsf{sound}}$ (negligible in the security parameter).
Therefore $P_{\mathcal{J}} \cap P_{\mathcal{E}} = \emptyset$.
\end{proof}

\subsection{Executive Containment}
\label{subsec:exec-containment}

\begin{theorem}[Executive Containment]
\label{thm:exec}
In a BTV institution $\mathcal{I}$, the operator $\mathcal{E}$ cannot
produce a decision $d$ without generating a non-repudiable evidence
receipt $e$ bound to $d$, and cannot delete any previously committed
receipt.
\end{theorem}

\begin{proof}
By the linear type invariant of Paper~1~\cite{btv2026},
$\mathtt{Verdict} \multimap E \otimes C$: a verdict value has
multiplicty 1 and can only be constructed by consuming exactly one
evidence token $E$ and one compliance certificate $C$. The linear
type system prevents duplication or discard of $E$: an unconsumed
evidence token is a compile-time error. Therefore every decision
generates exactly one receipt. By the append-only property of
Paper~2~\cite{btv-t-repo}, once a receipt is committed to $\Sigma$,
no capability in $P_{\mathcal{E}}$ supports a delete operation:
the Merkle tree admits only extension, never mutation. Together, these
two properties establish that $\mathcal{E}$ cannot produce a silent
decision (no receipt) or destroy an existing receipt. Executive
behavior is fully contained within the constitutional boundaries
defined by $\mathcal{L}$.
\end{proof}

\subsection{Completeness of the Separation}
\label{subsec:completeness}

\begin{corollary}[Constitutional Completeness]
\label{cor:completeness}
A BTV institution $\mathcal{I}$ satisfies the Constitutional Invariant
(Definition~\ref{def:invariant}): the capability sets of the three
layers are mutually disjoint.
\end{corollary}

\begin{proof}
Direct from Theorems~\ref{thm:leg},~\ref{thm:jud},
and~\ref{thm:exec}:
\begin{align*}
  P_{\mathcal{L}} \cap P_{\mathcal{E}} &= \emptyset
    &&\text{(Theorem~\ref{thm:leg})} \\
  P_{\mathcal{E}} \cap P_{\mathcal{J}} &= \emptyset
    &&\text{(Theorem~\ref{thm:jud})} \\
  P_{\mathcal{J}} \cap P_{\mathcal{L}} &= \emptyset
    &&\text{(Definition: $\mathcal{J}$ reads $\Sigma$; does not write
    $\mathcal{L}$)}
\end{align*}
Therefore $P_{\mathcal{L}} \cap P_{\mathcal{E}} \cap P_{\mathcal{J}}
= \emptyset$.
\end{proof}

\begin{remark}[Relation to TLA+ and Alloy]
Existing formal methods tools such as TLA+~\cite{lamport2002} and
Alloy~\cite{jackson2012} verify that a \emph{given specification}
satisfies logical invariants. Constitutional Code differs in scope:
it does not verify that a particular system is correct, but that the
\emph{architectural separation of capabilities} satisfies a political
theory requirement. TLA+ can verify Executive Containment
(Theorem~\ref{thm:exec}) for a specific operator implementation;
it cannot, without our framework, certify that the resulting institution
satisfies Legislative Independence or Judicial Independence as
political properties, because those concepts have no representation
in its specification language. Our contribution is the mapping between
political theory and cryptographic architecture, not the verification
tool.
\end{remark}
